\documentclass[10pt,a4paper]{amsart}
\usepackage[margin=19mm]{geometry}
\usepackage{amsmath,amssymb,mathtools}
\usepackage[scr=rsfs,cal=euler]{mathalfa}
\usepackage{xfrac}
\usepackage[unicode,hidelinks,bookmarks=false]{hyperref}
\newcommand{\ie}{i.e.\ }
\newcommand{\cf}{cf.\ }
\newcommand{\resp}{respectively\ }
\newcommand{\Subsection}{\subsection}
\providecommand{\nobreakdash}{\nobreak\hbox{-}}
\setlength{\emergencystretch}{2em}
\multlinegap0pt

\theoremstyle{plain}
\newtheorem{proposition}{Proposition}[section]
\newtheorem{conjecture}{Conjecture}[section]
\newtheorem{corollary}{Corollary}[section]
\newtheorem{lemma}{Lemma}[section]
\newtheorem{definition}{Definition}[section]
\newtheorem{theorem}{Theorem}[section]
\newtheorem{remark}{Remark}[section]
\newtheorem{example}{Example}[section]
\newtheorem{counterexample}{Counter Example}[section]
\newtheorem{observation}{Observation}[section]
\title[Distribution of the inversion statistic on run-sorted permutations]{Distribution of the inversion statistic on run-sorted permutations}
\author{Toufik Mansour, Olivia Nabawanda, Mark Shattuck}
\date{Source: arXiv:2609.28674v1 (CC BY 4.0)}
\begin{document}
\maketitle

\noindent\emph{Source and licence.} The delivered work is the article shown in the title above, version arXiv:2609.28674v1 (CC BY 4.0), distributed under the Creative Commons Attribution 4.0 International licence, as recorded in the arXiv licence metadata for this exact version and shown on the abstract page saved with this item. The full licence notice, which carries the licence URL and the address of that abstract page, is the bundled file \texttt{licenses/arxiv-2609.28674-CC-BY-4.0.txt}.
\par\smallskip

\noindent\textit{Editorial scope.} Counted unit: the article's Theorem (source label maxinvrnk) (source0.tex lines 587--598, source label \texttt{maxinvrnk}): ``If , then m n,k k 2 (r 2-r 2d-2)k 3 r 1 4 -2 r 1 3 (d-1) r 2 -dr d 2-5d 2 2 , where with and . If , then m n,k 3k 3r 2 4 k r 1 3 - 3r-2 4 r 1 3 - k r-d 1 3 - (rk k 2 d ) r 1 2 (k r) d 2 - d 3 , where ''. It is delivered together with the article's complete original proof (source0.tex lines 599--661, one proof environment adjacent to the statement, 63 lines), copied line for line from the article's own text. Required context retained uncounted: maxinv (lines 467--473). Editorial formatting only: an AMS \texttt{amsart} preamble that restates the article's own macro definitions and its own theorem declarations, and the theorem environments the retained lines use; the article's own class file is neither shipped nor input; the retained environments are renumbered sequentially in order of appearance, whereas the article prints them with its own numbering; the article's equation numbering is not reproduced and every cross-reference inside the retained text resolves to the wrapper's numbering. No citation occurs inside any retained line, so the wrapper carries no bibliography; All retained bodies are exact copies of the bundled \texttt{sources/211610/source0.tex}, so no omission receipt applies. No asset-inclusion command occurs inside any retained span and the package ships no image asset.


\section*{Retained context: maxinv (source0.tex lines 467--473)}

\begin{theorem}\label{maxinv}
Let $u_n$ denote the maximum number of inversions achieved by a member of $\mathcal{R}_n$, i.e., $u_n=\text{deg}(A_n(q,1))$.  Let $n=\binom{k}{2}+d$ for some $k \geq 2$ and $0 \leq d \leq k-1$.  Then we have
\begin{equation}\label{maxinve1}
u_n=\binom{k}{3}+3\binom{k}{4}+d\binom{k-1}{2}+\binom{d}{2},
\end{equation}
with this value being achieved by exactly $\binom{k-1}{d}+\binom{k-2}{d-2}$ members of $\mathcal{R}_n$.
\end{theorem}


\section{The counted unit: Theorem (source label maxinvrnk) and its complete original proof}

\begin{theorem}\label{maxinvrnk}
If $2k-1 \leq n <\binom{k+1}{2}$, then
\begin{equation}\label{maxinvrnke1}
m_{n,k}=k^2+(r^2-r+2d-2)k+3\binom{r+1}{4}-2\binom{r+1}{3}+(d-1)\binom{r}{2}-dr+\frac{d^2-5d+2}{2},
\end{equation}
where $n=2k-1+\binom{r}{2}+d$ with $1 \leq r \leq k-2$ and $0 \leq d <r$. If $n \geq \binom{k+1}{2}$, then
\begin{align}
m_{n,k}&=\frac{3k+3r+2}{4}\binom{k+r+1}{3}-\frac{3r-2}{4}\binom{r+1}{3}-\binom{k+r-d+1}{3}-\left(rk+\binom{k}{2}+d\right)\binom{r+1}{2}\notag\\
&\quad+(k+r)\binom{d}{2}-\binom{d}{3}, \label{maxinvrnke2}
\end{align}
where $n=\binom{k+1}{2}+rk+d$ with $r \geq 0$ and $0 \leq d <k$.  Furthermore, the value of $m_{n,k}$ is achieved by $\binom{r}{d}$ members of $\mathcal{R}_{n,k}$ in the first case and $\binom{k}{d}$ members of $\mathcal{R}_{n,k}$ in the second.
\end{theorem}

\begin{proof}
We modify the proof of Theorem \ref{maxinv}.  Since the procedures applied to non-maximal members of $\mathcal{R}_n$ described in that proof did not always preserve the number of runs, special care must be taken to obtain the maximal members of $\mathcal{R}_{n,k}$.  Let $\mathcal{R}_{n,k}'=\mathcal{R}_n'\cap \mathcal{R}_{n,k}$, where $\mathcal{R}_{n}'$ is as in the proof of Theorem \ref{maxinv}, and we may restrict attention to members of $\mathcal{R}_{n,k}'$.  As before, we represent $\pi \in \mathcal{R}_{n,k}'$ by a composition $a_\pi=(a_1,\ldots,a_k)$ wherein the $i$-th part of $a_\pi$ gives the length of the $i$-th run of $\pi$.  We first show \eqref{maxinvrnke1}, noting that $2k-1 \leq n <\binom{k+1}{2}$ for a fixed $k \geq 3$ may be uniquely represented by $k$, $r$ and $d$ in this case as indicated where $r$ and $d$ satisfy the stated restrictions.  If $r=1$, then $d=0$ and $n=2k-1$, with  $\mathcal{R}_{2k-1,k}'$ seen to consist of a single permutation $\pi$ whose associated composition is $(2,\ldots,2,1)$.  Note $\text{inv}(\pi)=\sum_{i=1}^{k-1}(2i-1)=(k-1)^2$, which agrees with the $r=1,\,d=0$ case of \eqref{maxinvrnke1}.

So let us assume $k \geq 4$, with $2 \leq r \leq k-2$, for the remainder of the proof of \eqref{maxinvrnke1}. Given $\pi \in \mathcal{R}_{n,k}'$ with $a_\pi=(a_1,\ldots,a_k)$, let $u_\pi=(u_1,\ldots,u_k)$ wherein $u_i=a_i-2$ for $1 \leq i < k$, with $u_k=a_k-1$, and let $v_\pi=(v_1,\ldots,v_k)$ wherein $v_i=u_i-(r-i)$ for $1 \leq i \leq r-1$, with $v_i=u_i$ for $r \leq i \leq k$.  Let $\mathcal{R}^{(k)}$ denote the subset of $\mathcal{R}_{n,k}'$ consisting of those members such that $v_i=0$ or $1$ for $1 \leq i \leq r$, with $v_{r+1}=\cdots=v_k=0$, and let $\mathcal{T}^{(k)}=\mathcal{R}_{n,k}'\backslash \mathcal{R}^{(k)}$.  We show that no member of $\mathcal{T}^{(k)}$ can achieve the maximum inv value of $m_{n,k}$. Let $\pi \in \mathcal{T}^{(k)}$ and first suppose $u_j \geq 1$ for some $j \in [r+1,k]$ within $u_\pi$.  Note we must have $v_i \leq 0$ for some $i \in [r-1]$, for otherwise
$\sum_{i=1}^ka_i \geq 2k-1+\binom{r}{2}+r>n$, which is impossible.  We then replace $a_j$ with $a_j-1$ and $a_i$ with $a_i+1$ within $a_\pi$ and consider the resulting member $\pi^* \in \mathcal{R}_{n,k}'$.

Proceeding as before, we have $\text{inv}(\pi^*)=\text{inv}(\pi)+a_j-a_i+j-i-1$.  Note $v_i \leq 0$ implies $a_i \leq r-i+2$, and hence
$$a_j-a_i+j-i-1 \geq a_j-(r-i+2)+j-i-1=a_j-3+j-r \geq 1,$$
as $u_j \geq 1$ implies $a_j \geq 3$ if $r<j<k$, with $a_k \geq 2$ if $j=k \geq r+2$.  Thus, $\text{inv}(\pi^*)>\text{inv}(\pi)$ and $\pi$ is not maximal, as desired.  So assume $u_j=0$ for each $j \in [r+1,k]$ within $u_\pi$.  Let $M=\max(v_i)_{1 \leq i \leq r}$ and $m=\min(v_i)_{1 \leq i \leq r}$.   Upon considering separately the cases (i) $M \geq 2$ and (ii) $M=1$ with $m<0$, one can show $\text{inv}(\pi)<m_{n,k}$ for the remaining $\pi \in \mathcal{T}^{(k)}$, by applying an argument comparable to the one used at a similar juncture in the proof of Theorem \ref{maxinv}.  Note $M \geq 2$ in (i) implies $m \leq 0$, for otherwise the sum of the components of $a_\pi$ would exceed $n$.

One may verify that all members of $\mathcal{R}^{(k)}$ have equal inv value, with $|\mathcal{R}^{(k)}|=\binom{r}{d}$.  We first compute $\text{inv}(\sigma)$ when $d=0$, where $\sigma$ denotes the sole member of $\mathcal{R}^{(k)}$ in this case.  Note $\sigma \in \mathcal{R}^{(k)}$ together with  $n=2k-1+\binom{r}{2}$  implies $u_\sigma=(r-1,r-2\ldots,1,0,\ldots,0)$.  Upon considering separately the inversions involving the second elements in each of the runs of $\sigma$, we have
\begin{align*}
\text{inv}(\sigma)&=\sum_{i=1}^{k-1}(2i-1)+\sum_{i=0}^{r-3}(r-i)\binom{r-i-1}{2}+\sum_{i=1}^{r-1}(r-i)(2(k-i)-1)\\
&=(k-1)^2+\sum_{i=1}^{r-1}(i+1)\binom{i}{2}+\sum_{i=1}^{r-1}i(2(k-r+i)-1)\\
&=(k-1)^2+\frac{3}{2}\sum_{i=1}^{r-1}i^2+\frac{1}{2}\sum_{i=1}^{r-1}i^3 +(2k-2r-1)\binom{r}{2}+\binom{r}{3}\\
&=(k-1)^2+\frac{r(r-1)(2r-1)}{4}+\frac{1}{2}\binom{r}{2}^2+(2k-2r-1)\binom{r}{2}+\binom{r}{3}.
\end{align*}
This last expression may be rewritten somewhat to give
\begin{equation}\label{d=0case}
\text{inv}(\sigma)=(k-1)^2+kr(r-1)+3\binom{r+1}{4}-2\binom{r+1}{3}-\binom{r}{2},
\end{equation}
which yields the $d=0$ case of \eqref{maxinvrnke1}.

Now suppose $\tau \in \mathcal{R}^{(k)}$, where $d \in [r-1]$.  We may assume $a_\tau$ is obtained from $a_\sigma$ be adding 1 to each of the first $d$ entries of $a_\sigma$. Considering separately from the others those inversions that arise due to this addition of 1 to the first $d$ entries of $a_\sigma$, we have
\begin{align*}
\text{inv}(\tau)&=\text{inv}(\sigma)+\binom{d}{2}+\sum_{i=1}^d\binom{r-i}{2}+\sum_{i=1}^d(2k-2i-1)+\sum_{i=0}^{d-2}(r-i)(d-i-1)\\
&=\text{inv}(\sigma)+\binom{d}{2}+\binom{r}{3}-\binom{r-d}{3}+d(2k-1)-d(d+1)+\sum_{i=1}^{d-1}i(r-d+i+1)\\
&=\text{inv}(\sigma)+\binom{d}{2}+\binom{r}{3}-\binom{r-d}{3}+d(2k-1)-d(d+1)+(r-d)\binom{d}{2}+2\binom{d+1}{3}\\
&=\text{inv}(\sigma)+2dk+\binom{r}{3}-\binom{r-d}{3}+r\binom{d}{2}-\frac{d(d^2+17)}{6}.
\end{align*}
Using formula \eqref{d=0case} for $\text{inv}(\sigma)$, and combining some terms, we get
$$\text{inv}(\tau)=k^2+(r^2-r+2d-2)k+3\binom{r+1}{4}-\binom{r+1}{3}-2\binom{r}{2}-\binom{r-d}{3}+r\binom{d}{2}-\frac{d^3+17d-6}{6}.$$
Applying the fact
$$r\binom{d}{2}-\binom{r-d}{3}=d\binom{r}{2}+\binom{d+2}{3}-\binom{r}{3}-dr$$
to the last expression yields \eqref{maxinvrnke1} after some algebraic steps.

We now show \eqref{maxinvrnke2}. First note $n \geq \binom{k+1}{2}$ may be uniquely represented in terms of $k$, $r$ and $d$ as described.  Further, formula \eqref{maxinvrnke2} is seen to give the correct value of zero for all $n \geq 1$ when $k=1$, so we may assume $k \geq 2$.  Let $\pi \in \mathcal{R}_{n,k}'$, where $n=\binom{k+1}{2}+rk+d$ and $r$ and $d$ are as stated.  If $a_\pi=(a_1,\ldots,a_k)$, then define the vector $w_\pi=(w_1,\ldots,w_k)$ wherein $w_i=a_i-k-r+i-1$ for $1 \leq i \leq k$. Let $\mathcal{R}^{(k)}$ in this case denote the subset of $\mathcal{R}_{n,k}'$ consisting of those $\pi$ such that $w_\pi$ consists of only $0$'s and $1$'s and let $\mathcal{T}^{(k)}=\mathcal{R}_{n,k}'\backslash \mathcal{R}^{(k)}$.  Note $d<k$ in the representation of $n$ implies $w_\pi$ cannot be all $1$'s. Let $M=\max(w_i)_{1 \leq i \leq k}$ and $m=\min(w_i)_{1 \leq i \leq k}$, and note $M \geq 1$ with $M -m \geq 2$ for all members of $\mathcal{T}^{(k)}$.  Proceeding as before considering $M$ and $m$, one can show that no member of $\mathcal{T}^{(k)}$ can achieve the maximum inv value on $\mathcal{R}_{n,k}$.  Further, it is seen that all members of $\mathcal{R}^{(k)}$ achieve the same (maximum) inv value, with $|\mathcal{R}^{(k)}|=\binom{k}{d}$.

Let $\sigma$ denote the sole member of $\mathcal{R}^{(k)}$, where $n$ is of the stated form with $d=0$, and let $\tau$ be a member of $\mathcal{R}^{(k)}$, where $d \in [k-1]$ in the representation of $n$.  Note $a_\sigma=(k+r,k+r-1,\ldots,r+1)$ and we may assume $\tau$ arises by adding $1$ to each of the first $d$ entries of $a_\sigma$ and considering the resulting permutation.  We first compute $\text{inv}(\sigma)$.  Upon considering the number of inversions caused by the letters in the $(k-j-1)$-st run of $\sigma$ for $0 \leq j \leq k-2$, we have
\begin{align*}
\text{inv}(\sigma)&=\sum_{j=0}^{k-2}(j+r+1)\sum_{i=r+1}^{j+r+1}i=\sum_{j=0}^{k-2}(j+r+1)\left(\binom{j+r+2}{2}-\binom{r+1}{2}\right)\\
&=\sum_{j=0}^{k-2}(j+r+1)\binom{j+r+2}{2}-\binom{r+1}{2}\sum_{j=0}^{k-2}(j+r+1)\\
&=3\sum_{j=0}^{k-2}\binom{j+r+3}{3}-2\sum_{j=0}^{k-2}\binom{j+r+2}{2}-\left(r(k-1)+\binom{k}{2}\right)\binom{r+1}{2}\\
&=3\binom{k+r+2}{4}-3\binom{r+3}{4}-2\binom{k+r+1}{3}+2\binom{r+2}{3}-\left(r(k-1)+\binom{k}{2}\right)\binom{r+1}{2}\\
&=\binom{k+r+2}{4}+2\binom{k+r+1}{4}-\binom{r+3}{4}-2\binom{r+2}{4}-\left(r(k-1)+\binom{k}{2}\right)\binom{r+1}{2}.
\end{align*}
Using the fact $\binom{m}{4}+2\binom{m-1}{4}=\frac{3m-8}{4}\binom{m-1}{3}$ for $m \geq 3$ in the last expression implies
\begin{align}
\text{inv}(\sigma)&=\frac{3k+3r-2}{4}\binom{k+r+1}{3}-\frac{3r+1}{4}\binom{r+2}{3}-\left(r(k-1)+\binom{k}{2}\right)\binom{r+1}{2}\notag\\
&=\frac{3k+3r-2}{4}\binom{k+r+1}{3}-\frac{3r-2}{4}\binom{r+1}{3}-\left(rk+\binom{k}{2}\right)\binom{r+1}{2},\label{d=0case(b)}
\end{align}
which establishes the $d=0$ case of \eqref{maxinvrnke2}.

Considering the inversions caused by the extra elements added to the first $d$ runs of $\sigma$ in obtaining $\tau$, we have
$$\text{inv}(\tau)=\text{inv}(\sigma)+\binom{d}{2}+\sum_{j=1}^d\sum_{i=r+1}^{k+r-j}i+\sum_{i=0}^{d-2}(d-i-1)(k+r-i-1).$$
Note
$$\sum_{j=1}^d\sum_{i=r+1}^{k+r-j}i=\sum_{j=1}^d\left(\binom{k+r-j+1}{2}-\binom{r+1}{2}\right)=\binom{k+r+1}{3}-\binom{k+r-d+1}{3}-d\binom{r+1}{2}$$
and
$$\sum_{i=0}^{d-2}(d-i-1)(k+r-i-1)=(k+r-d)\binom{d}{2}+\frac{d(d-1)(2d-1)}{6}=(k+r)\binom{d}{2}-\binom{d+1}{3},$$
which implies
$$\text{inv}(\tau)=\text{inv}(\sigma)+\binom{k+r+1}{3}-\binom{k+r-d+1}{3}-d\binom{r+1}{2}+(k+r)\binom{d}{2}-\binom{d}{3}.$$
Formula \eqref{maxinvrnke2} now follows from \eqref{d=0case(b)}, which completes the proof.
\end{proof}

\end{document}
